\documentclass[12pt]{article}
\usepackage[utf8]{inputenc}
\usepackage{amsmath}
\usepackage{amssymb}
\usepackage{graphicx}
\usepackage{booktabs}
\usepackage{listings}
\usepackage{xcolor}
\usepackage{hyperref}
\usepackage{geometry}
\geometry{margin=1in}

\lstset{
    language=Python,
    basicstyle=\ttfamily\small,
    keywordstyle=\color{blue}\bfseries,
    commentstyle=\color{green!60!black},
    stringstyle=\color{red},
    numbers=left,
    numberstyle=\tiny\color{gray},
    stepnumber=1,
    numbersep=5pt,
    backgroundcolor=\color{gray!10},
    frame=single,
    breaklines=true,
    breakatwhitespace=false,
    tabsize=4,
    showstringspaces=false
}

\title{Build Costs Calculation Documentation}
\author{}
\date{}

\begin{document}

\maketitle

\section{Introduction}

The \texttt{build\_costs.py} script calculates the construction costs for a transmission line project. It computes costs for conductors, structures, and converters, applies terrain adjustments via weighted miles, incorporates contingency factors, and calculates both regulatory (AFUDC) and societal (present value) perspectives.

\section{Method Overview}

\subsection{Purpose}

This script calculates the total build cost of a transmission line project by:
\begin{enumerate}
    \item Loading base costs from YAML configuration files
    \item Adjusting costs for terrain complexity using weighted miles
    \item Applying contingency factors to account for uncertainty
    \item Calculating AFUDC (Allowance for Funds Used During Construction) for regulatory perspective
    \item Computing present value and amortized costs for societal perspective
\end{enumerate}

\subsection{Calculation Approach}

The build cost calculation follows this structure:

\begin{align}
\text{Conductor Cost} &= (\text{Variable Cost/Mile} \times \text{Weighted Miles}) + \text{Fixed Cost} \\
\text{Structure Cost} &= \text{Variable Cost/Mile} \times \text{Weighted Miles} \quad \text{(if not reconductoring)} \\
\text{Converter Cost} &= \text{Fixed Cost} \times \text{Number of Converters} \quad \text{(if not reconductoring)} \\
\text{Total Base Cost} &= \text{Conductor} + \text{Structure} + \text{Converter} \\
\text{Total with Contingencies} &= \text{Base Cost} \times (1 + \text{Contingency Rate})
\end{align}

For reconductoring projects, structure and converter costs are set to zero since existing infrastructure is reused.

\subsection{Inputs and Outputs}

\textbf{Inputs:}
\begin{itemize}
    \item Project technical details (construction type, AC/DC, capacity, conductor type, etc.)
    \item Physical details (terrain distribution)
    \item Build cost parameters from YAML files
    \item Contingency percentages
    \item Financing parameters (WACC, inflation, AFUDC rate)
\end{itemize}

\textbf{Outputs:}
\begin{itemize}
    \item Nominal costs (with contingencies)
    \item AFUDC capitalized costs (regulatory perspective)
    \item Present value costs (societal perspective)
    \item Component-level costs (conductor, structure, converter)
\end{itemize}

\subsection{Integration with Other Scripts}

This script integrates with:
\begin{itemize}
    \item \texttt{weighted\_miles.py} - Calculates terrain-adjusted miles (see \texttt{utilities\_documentation.tex})
    \item \texttt{yaml\_loaders.py} - Loads project configuration (see \texttt{utilities\_documentation.tex})
    \item \texttt{financial\_utils.py} - Financial calculations (see \texttt{utilities\_documentation.tex})
    \item \texttt{csv\_output\_manager.py} - Writes results to CSV (see \texttt{csv\_output\_manager\_documentation.tex})
\end{itemize}

\section{Key Functions}

\subsection{\texttt{load\_costs(category, total\_miles, number\_of\_converters, contingencies, reconductoring)}}

Loads and calculates build costs for the specified project category.

\textbf{Parameters:}
\begin{itemize}
    \item \texttt{category} (str): Project category identifier in format \texttt{"ConstructionType/AC\_DC/CapacityMW/ConductorType/ConverterType"}
    \item \texttt{total\_miles} (float): Total miles of transmission line
    \item \texttt{number\_of\_converters} (int): Number of converters needed (0 for AC)
    \item \texttt{contingencies} (dict): Dictionary with keys \texttt{"conductor\_contingency"}, \texttt{"structure\_contingency"}, \texttt{"converter\_contingency"}
    \item \texttt{reconductoring} (bool): Whether this is a reconductoring project
\end{itemize}

\textbf{Returns:}
\begin{itemize}
    \item \texttt{total\_cost} (float): Base total cost without contingencies
    \item \texttt{total\_cost\_with\_contingencies} (float): Total cost with contingencies applied
    \item \texttt{conductor\_cost} (float): Base conductor cost
    \item \texttt{structure\_cost} (float): Base structure cost (0 for reconductoring)
    \item \texttt{converter\_cost} (float): Base converter cost (0 for reconductoring)
    \item \texttt{conductor\_cost\_with\_contingencies} (float): Conductor cost with contingency
    \item \texttt{structure\_cost\_with\_contingencies} (float): Structure cost with contingency
    \item \texttt{converter\_cost\_with\_contingencies} (float): Converter cost with contingency
    \item \texttt{weighted\_miles} (float): Terrain-adjusted miles
    \item \texttt{average\_terrain\_multiplier} (float): Average terrain cost multiplier
\end{itemize}

\textbf{Key Logic:}
\begin{lstlisting}
# Calculate weighted miles for terrain adjustment
weighted_miles, average_terrain_multiplier = calculate_weighted_miles()

# Conductor cost: variable per mile + fixed
conductor_cost = (variable_conductor_cost_per_mile * weighted_miles) 
                 + fixed_conductor_cost

# Structure and converter costs only for new builds
if reconductoring:
    structure_cost = 0
    converter_cost = 0
else:
    structure_cost = variable_structure_cost_per_mile * weighted_miles
    converter_cost = fixed_converter_cost * number_of_converters

# Apply contingencies
conductor_cost_with_contingencies = conductor_cost * 
    (1 + contingencies["conductor_contingency"])
\end{lstlisting}

\subsection{\texttt{main()}}

Main execution function that orchestrates the build cost calculation and output.

\textbf{Workflow:}
\begin{enumerate}
    \item Loads project technical details and constructs category identifier
    \item Determines number of converters (0 for AC projects)
    \item Loads physical details, contingencies, and financing parameters
    \item Calls \texttt{load\_costs()} to calculate base costs
    \item Calculates AFUDC capitalized costs (if enabled)
    \item Calculates amortized costs and present value
    \item Displays results and writes to CSV
\end{enumerate}

\section{Example}

The following example demonstrates a typical usage of the build costs calculation:

\begin{lstlisting}
# The script is typically called from ctcc.py, but can be run standalone
# It automatically loads configuration from YAML files

# Example: 500 MW Overhead AC line with Advanced Aluminum Conductor
# Category: "Overhead/AC/500MW/Advanced Aluminum/NA"
# 
# Inputs (from YAML):
#   - Construction type: Overhead
#   - AC/DC: AC
#   - Capacity: 500 MW
#   - Conductor: Advanced Aluminum
#   - Total miles: 100 miles
#   - Terrain: 60% flat, 40% forested
#   - Contingencies: 15% conductor, 20% structure, 10% converter
#
# Calculation:
#   1. Weighted miles = 100 * (0.6*1.0 + 0.4*1.2) = 108 miles
#   2. Conductor cost = ($500k/mile * 108) + $1M = $55M
#   3. Structure cost = $200k/mile * 108 = $21.6M
#   4. Converter cost = 0 (AC project)
#   5. Total base = $76.6M
#   6. With contingencies = $76.6M * (1 + weighted avg) ≈ $88M
#
# Output:
#   - Nominal cost: $88M
#   - AFUDC capitalized: $92M (if AFUDC enabled)
#   - Present value: $88M (discounted at real WACC)
\end{lstlisting}

\textbf{Integration Example:}

\begin{lstlisting}
# When called from ctcc.py:
# 1. Script loads project details from YAML
# 2. Calculates build costs
# 3. Writes to CSV via CTCCOutputManager:
csv_manager = CTCCOutputManager()
results = {
    "total_nominal": total_cost_with_contingencies,
    "total_afudc": capitalized_cost_with_contingencies,
    "total_pv": pv_of_amortized,
    # ... component costs
}
csv_manager.add_build_costs(results)
csv_manager.write_batch_summary()
\end{lstlisting}

\section{Key Implementation Details}

\subsection{Reconductoring Handling}

For reconductoring projects, the script sets structure and converter costs to zero since existing infrastructure is reused. Only conductor costs are calculated.

\subsection{AFUDC Calculation}

If AFUDC is enabled, the script calculates capitalized costs using the AFUDC rate and cost timing patterns. This represents the regulatory perspective where construction financing costs are capitalized.

\subsection{Present Value Verification}

The script verifies amortization calculations by discounting annual amortized payments back to present value, which should equal the original cost (within rounding).

\subsection{Weighted Miles}

The script uses \texttt{calculate\_weighted\_miles()} to adjust costs for terrain complexity. Different terrain types have different cost multipliers, and the weighted average is applied to the total miles.

\end{document}

